\section{Conclusion}
\label{sec:conclusion}

We set out to check a 12.74$\times$ speedup for parallel analysis in a selective
defragmenter and found that the speedup belonged to the delay model, not to the
storage. The parallel mode overlapped sleeps placed before an extent lookup that the
original analysis never needed. The original's real costs were process creation and
repeated tree walks, which a short single-threaded rewrite removes with
byte-identical output. Delays injected below the cache, where real devices impose
them, leave threads little to overlap when files are created together and only a
bounded gain when they are scattered. For tools that do per-file metadata work, the
habits that would have caught this are cheap: a strong sequential baseline, latency
injected at the block layer, and a count of device reads per unit of work.
